\Qhead{4.7}{Усилители СВЧ-диапазона}{Усилители СВЧ-диапазона (резонаторный, бегущей волны). Полоса пропускания усилителя бегущей волны.}
\Subfile{4.7.1}{sections/04_generation/q07/s01}

\recap{Усилитель извлекает энергию постоянного питания через фазово согласованный переменный ток. Резонаторный вариант ограничен полосой моды, бегущая волна — рассогласованием пучка и дисперсионной ветви.}
\board{Схема вход–замедляющая система–коллектор; диаграмма $\omega(k)$; накопленная фазовая ошибка.}
\oral{Почему увеличение длины взаимодействия может повышать усиление, но сужать полезную полосу?}
\bait{Можно оставить различие фазовой и групповой скорости: точка синхронизма задаётся первой, изменение расстройки с частотой — второй.}
\sources{Carter~\cite{carter}, Radio-frequency power generation, разделы об электронных усилителях; Белокопытов~\cite{belokopytov1}, связанные волны. Пособие Саввина из программы полностью не получено; конкретная формула полосы ему не приписывается.}
